\documentclass{article}
\usepackage[utf8]{inputenc}
\usepackage{xspace}
\usepackage[top=1.5cm, bottom=1.5cm, left=1.5cm, right=1.5cm]{geometry}
\usepackage{enumitem}
\usepackage{graphicx}
\usepackage{textcomp}
\usepackage{amssymb}
\usepackage{ifsym}
\usepackage{xspace}
\usepackage{amsmath}
\usepackage{stmaryrd}

\title{Facultad de Ciencias. \\Tarea 5 \\ Fundamentos de Bases de             Datos.}
\author{Arteaga Hernández Alejandro Iabin \and 
        Barocio Gómez Luis Fernando \and 
        Ledesma Granados Roberto Andrés \and
        Castrejón Estrada Juan Carlos}
\date{Entrega 4 de Mayo 2020}

\begin{document}

\maketitle
\begin{enumerate}
    %%%Pregunta 1%%%
    \item Preguntas de repaso:
    
    \begin{itemize}
    
        \item Anomalía de inserción: Cuando deseamos agregar datos a una base de datos, pero es imposible hacerlo, debido a la ausencia de otros datos.\\
        Anomalía de borrado: Se eliminan datos de forma no intencional, debido a que se eliminarion otros datos.
        \\ Anomalía de actualización: Se presentan incosistencias en los datos, debido a actualizaciones parciales.
        \item Las dependencias funcionales nos ayudan a definir cuando un diseño de bases de datos es correcto.
        \\ Es una relación entre dos atributos A y B, de forma que para cada valor único de A, sólo un valor de B se asocia con él a través de ésta relación.
        \item Uno de los objetivos de la normalización es conservar las dependencias funcionales. Ésto puede servir para acabar con las anomalías de redundancia.
        \item $A \rightarrow A, A \rightarrow B, A \rightarrow C, A \rightarrow AB, A \rightarrow AC, A \rightarrow BC, A \rightarrow ABC.$
        \item La normalización es una técnica para diseñar la estructura lógica de una base de datos en el modelo relacional. Para llegar a ésto, existe un proceso en el cual se van comprobando el cumplimiento de una serie de restricciones. Una forma normal consiste en un conjunto de éstas restricciones.
        \item Es preferible la 3NF cuando no queremos perder nada de información sobre la tabla orignial.
    \end{itemize}
     %%%Pregunta 2%%%
    \item Proporciona algunos ejemplos que demuestren que las siguientes reglas no son válidas:
    \begin{enumerate}
        \item Si A $\rightarrow$ B, entonces B $\rightarrow$ A.\\
        Consideremos la siguiente tabla:\\
        A    B\\
        a1   b1\\
        a2   b2\\
        a3   b2\\
        Claramente A $\rightarrow$ B, pues las a's son únicas, pero no es cierto que B $\rightarrow$ A.
        
        \item Si AB $\rightarrow$ C, entonces A $\rightarrow$ C y B $\rightarrow$ C.\\
        Consideremos la siguiente tabla:\\
        A     B     C\\
        a1    b1    c1\\
        a2    b1    c2\\
        Tenemos que AB $\rightarrow$ C, pero no es cierto que B $\rightarrow$ C.
        \item Si AB $\rightarrow$ C y A $\rightarrow$ C, entonces B $\rightarrow$ C.\\
        Consideremos la table del ejemplo anterior. En ella, AB $\rightarrow$ C y A $\rightarrow$ C, pero no es cierto que B $\rightarrow$ C.
    \end{enumerate}

     %%%Pregunta 3%%%
    \item Dada una relación $R (A, B, C, D, E, G)$ y el siguiente conjunto de dependencias funcionales $F=${ $AB\rightarrow C, BC \rightarrow D,$\\$ D \rightarrow EG, CG \rightarrow BD, C \rightarrow A, ACD \rightarrow B, BE \rightarrow C, CE \rightarrow AG$}.
    \begin{enumerate}
        \item La cerradura de BC es {A, D, E, G}.\\
        Es falso. En realidad, la cerradura de BC es ABCDEG.\\
        Por reflexividad y descomposición, tenemos que BC $\rightarrow$ B y C $\rightarrow$ C. Tenemos de forma directa que BC $\rightarrow$ D, y por reflexividad con C tenemos que BC $\rightarrow$ A. Por reflexividad con D y descomposición, tenemos que BC $\rightarrow$ E y BC $\rightarrow$ G.
        \item Todos los atributos de R están en la cerradura de BC.\\
        Es cierto. El argumento es el mismo que en el ejercicio anterior.
        \item La cerradura de AC es {A, C}. Es cierto. No podemos obtener más elementos en la cerradura más que A y C.
        \item ABC es una superllave de R.
        \\ Es cierto. Dado que BC es una llave (justificado en el ejercicio (a)), ABC es una superllave de R por contener una llave.
        \item ABC es una llave candidata de R.
        \\Esto es falso. ABC no es una superllave mínima, ya que BC es llave y por lo tanto, superllave.
        \item BC es la única llave candidata de R.\\
        Esto es falso. Notemos que BE $\rightarrow$ C, y BE $\rightarrow$ B, por lo que BE $\rightarrow$ BC. Y como BC era llave, BE también debe serlo, por lo que BE es otra llave candidata.
    \end{enumerate}
    
    %%%Pregunta 4%%%
    \item Sean $R (A, B, C, D, E, F, G)$, el conjunto $F = \{ BD \rightarrow A, C \rightarrow D, G \rightarrow C, E \rightarrow FG\}$. Especifica cuáles de las siguientes tuplas podrían estar almacenadas en R y por qué:
    \begin{itemize}
        \item (a1, b1, c1, d1, e2, f1, g2): No la podemos agregar. Si lo hacemos, tendríamos dos valores iguales en C (c1) que definen dos valores distintos de D (d1 y d2).
        \item (a1, b3, c2, d2, e1, f1, g1): No la podemos agregar. Si lo hacemos, tendríamos dos valores iguales en G (g1) que definen dos valores distintos de C (c1 y c2).
        \item (a1, b2, c3, d1, e4, f2, g3): Si la podemos agregar. No causa conflicto con ninguna de las dependencias funcionales y la tabla R.
        \item (a1, b1, c2, d2, e2, f1, g2): Si la podemos agregar. No causa conflicto con ninguna de las dependencias funcionales y la tabla R.
    \end{itemize}
    
    %%%Pregunta 5%%%
    \item Dado el conjunto de dependencias funcionales $F = \{ A \rightarrow DC, BC \rightarrow F, F \rightarrow CH, AF \rightarrow B \} $
    \begin{itemize}
        \item Indica alguna llave candidata 
        
        $\{A,F\}$
        
        \item Encontrar las cerraduras: 
        
        $\{AC\}+ = \{ ACD\}$,
        
        $\{BA\}+ = \{ ABCDFH\}$, 
        
        $\{AF\}+ = \{ ABCDFH\}$ y 
        
        $ \{ A \}+ = \{ ADC\}$
        
        \item Indicar si las siguientes dependencias funcionales son deducidas o no, justificando tu respuesta:
        
        $ AC \rightarrow H$, No se puede porque H solo depende de F y F depende de B y B no
        depende de AC
        
        $BA \rightarrow FC$, si $BA \rightarrow BC->FC$
        
        $AF \rightarrow HB$  Si porque $ AF \rightarrow B $ y  $ F \rightarrow CH$
        
        $A \rightarrow F$ NO porque A no deriva a BC y es la que deriva a F
        
    \end{itemize}
    
    
     %%%Pregunta 6%%%
    \item Para cada uno de los esquemas que se muestran a continuación:
    \begin{enumerate}
        \item $R (A, B, C, D, E)$ con $F = {AB \rightarrow CD, E \rightarrow C, D \rightarrow B}$
        
        \item $R (A, B, C, D, E)$ con $F = {AB \rightarrow C, DE \rightarrow C, B \rightarrow D}$
        
        \begin{itemize}
            \item Especifica de ser posible \textbf{dos DF no triviales} que se pueden derivar de las dependencias funcionales dadas
            
            a) $AD \rightarrow C$, $AB\rightarrow C$
            
            b) $AB \rightarrow CD$, $BE \rightarrow C$
            
            \item Indica \textbf{alguna llave candidata} para \textbf{R}
            $\{E A B\}$
            \item Especifica \textbf{todas las violaciones} a la \textbf{BCNF}
            
            a) Ninguna DF cumple BCNF porque ningún determinante de ninguna Df es superllave
            
            b) Ninguna DF cumple BCNF porque ningún determinante de ninguna Df es superllave
            
            \item \textbf{Normaliza} de acuerdo a \textbf{BCNF}, asegúrate de indicar cuáles son las relaciones resultantes con sus respectivas dependencias funcionales.
            
            \begin{itemize}
                \item $AD \rightarrow C$, $AB\rightarrow C$
                
                $\{AB \}+=\{ABCD\}\times$(Llave por R = ABE)
                
                $\{E \}+= \{EC\} \times$
                
                $\{D \}+=\{DB\} \times $(Llave por R = ABE)
                
                \textbf{S(A,B,C,D)} con $\{AB \rightarrow CD\,D \rightarrow B \}$
                
                $\{AB\}+=\{ABCD\}$ por lo tanto AB = llave
                
                $\{D \}+={DB} \times$
                
                $u=(D,B)$ con $D\rightarrow B$ por lo tanto D = llave, cumple. 
                
                $V(D,A,C)$ con $DAC \rightarrow DAC$ por lo tanto $DAC = llave$ tambien cumple.
                
                DF es perdida \{$E \rightarrow C$ y $AB \rightarrow CD$ \}(Join con perdida)
                
                Por lo tanto BCNF no sirve y esto es 3NF
                
                \textbf{T(A,B,E)} con $ABE \rightarrow ABE$ por lo tanto $ABE = llave$
                
              
            \end{itemize}
            
            
        \end{itemize}
        
    \end{enumerate}
    
    %%%Pregunta 7%%%
    \item Para cada una de las siguientes relaciones con su respectivo conjunto de dependencias funcionales:
    \begin{enumerate}
        \item $R (A, B, C, D, E, F)$ con $F = \{ B \rightarrow D, B \rightarrow E, D \rightarrow F, AB \rightarrow C \}$
        
        Calcular conjunto F mínimo:
        
        Como $F$ tiene $B \rightarrow D$ y $B \rightarrow E$, aplicamos la regla de la unión y sustituimos estas $DP$ con $B \rightarrow DE$, lo cual nos da $F = \{ B \rightarrow DE, D \rightarrow F, AB \rightarrow C \}$
        
        Ahora eliminaremos los atributos superfluos (por la izquierda) de las DP:
        
        $AB \rightarrow C \rightarrowtriangle A \rightarrow C:$
        
        $\{A\}+ = \{AC\}, B$ no es superfluo
        
        $AB \rightarrow C \rightarrowtriangle B \rightarrow C:$
        
        $\{B\}+ = \{BDEFC\}, A$ no es superfluo
        
        No hay cambios.
        
        Ahora por la derecha:
        
        $B \rightarrow DE \rightarrowtriangle B \rightarrow D:$
        
        $\{B\}+ = \{BDF\}, E$ no es superfluo
        
        $B \rightarrow DE \rightarrowtriangle B \rightarrow E:$
        
        $\{B\}+ = \{BE\}, D$ no es superfluo
        
        No hay cambios.

        $F = \{ B \rightarrow DE, D \rightarrow F, AB \rightarrow C\}$ es un conjunto mínimo.
        
        $S(B, D, E), T(D, F), U(A, B, C)$ no son subconjunto entre ellos.
        
        La llave candidata $AB (\{AB\}+ = \{ABCDEF\})$ se encuentra en $U$, por lo que terminamos con la normalización, y nos queda:
        
        $S(B, D, E), T(D, F), U(A, B, C)$

        
        \item $R (A, B, C, D, E, G, H)$ con $F = \{C \rightarrow D, A \rightarrow C, AD \rightarrow GBH, B \rightarrow A, BD \rightarrow H\}$ 
        
    \end{enumerate}
    

     \begin{itemize}
         \item Indica \textbf{alguna llave} para la relación R e indica las \textbf{violaciones} a \textbf{3NF} que encuentres en F
         
            $\{A\}$ es llave primaria.
            
            $C \rightarrow D$ Viola 3FN porque no está la llave A,
            De igual manera $ BD \rightarrow H$
         
         \item Encuentra el \textbf{conjunto mínimo} de dependencias funcionales equivalente a F.
        se calcularon los atributos superfluos por la izquierda 
        
         $ C \rightarrow D$
         
         $A \rightarrow C$
         
         $A \rightarrow GBH $
         
         $B \rightarrow A$
         
         $B \rightarrow H $
         
         ahora por la derecha
         
         $ C \rightarrow D$
         
         $A \rightarrow C$
         
         $A \rightarrow GB$
         
         $B \rightarrow A$
         
         $B \rightarrow H $

         
         \item \textbf{Normaliza} de acuerdo a la \textbf{3NF}. Indica claramente las relaciones resultantes y en cada esquema, las dependencias funcionales que se cumplen.

         
        R1(C, D) $ C \rightarrow D$
         
        R2(A, C) $A \rightarrow C$
         
        R3(A, G, B) $A \rightarrow GB $
         
        R4(B, A) $B \rightarrow A $
        
        R5(B, H) $B \rightarrow H $
         
        
    \end{itemize}
    
    
    %%%Pregunta 8%%%
    \item Sea el esquema:
    
    $R (A, B, C, D, E, F)$ con $F = \{ BD \rightarrow E, CD \rightarrow A, E \rightarrow C, B \rightarrow D \}$
    
    \begin{itemize}
        \item ¿Qué puedes decir de $\{A\}+$ y $\{F\}+$?
        
        $\{A\}+ =  \{A\}$ y $\{F\}+ = \{F\}$ Porque ninguna entidad depende funcionalmente
        de ellas.
        
        \item Calcula $\{B\}+$, ¿qué puedes decir de esta cerradura?

        $\{B\}+ =  \{A, B, C, D, E\}$ No es una llave porque F no depende funcionalmente  
        
        \item Obtén todas las \textbf{llaves candidatas}.
        
        $\{B, F \}$ es una llave mínima porque son los únicos atributos que no dependen, funcionalmente de ningún otro.
        
        \item ¿R cumple con \textbf{BCNF}? ¿Cumple con \textbf{3NF}? (en caso contrario normaliza)
        
        No cumple ninguna de las 2, porque la llave no aparece en ninguna dependencia funcional
        
        Normalizaremos a 3FN, primero encontraremos el junto mínimo.
        $B  \rightarrow E$
        
        $CD  \rightarrow A$
        
        $E  \rightarrow C$
        
        $B  \rightarrow D$
        
        Crearemos una tabla para cada df minimizada
        
        R1(B, E) $B  \rightarrow E$
        
        R2(C, D, A) $CD  \rightarrow A$
        
        R3(E, C) $E  \rightarrow C$
        
        R4(B, D) $B  \rightarrow D$
        
        Y como la llave no está en ninguna de las relaciones se agrega la relación
        
        R5(B, F)
        
        \item Se ha decidido dividir \textbf{R} en las siguientes relaciones $S (A, B, C, D, F)$ y $T (C, E)$, ¿se puede recuperar la información de \textbf{R}?
        
        No, la dependencia funcional $ BD \rightarrow E$ se perdería.

    \end{itemize}
    
\end{enumerate}

\end{document}
